% Student-facing guide to the two-pass AI Fluency quiz format.
% Compile with: latexmk -pdf ai-fluency-quiz-student-guide.tex
\documentclass[10pt]{article}

\usepackage[letterpaper,left=0.82in,right=0.82in,top=0.72in,bottom=0.72in,headheight=13pt,footskip=22pt]{geometry}
\usepackage[T1]{fontenc}
\usepackage{lmodern}
\usepackage{microtype}
\usepackage{array,booktabs,tabularx}
\usepackage{enumitem}
\usepackage{fancyhdr}
\usepackage[hidelinks]{hyperref}

\newcommand{\CourseCodes}{ECE 4424 / CS 4824}
\newcommand{\CourseName}{Machine Learning: Learn by Building}
\newcommand{\CourseTerm}{Fall 2026}
\newcommand{\InstructorName}{Ming Jin}

\setlength{\parindent}{0pt}
\setlength{\parskip}{0.45em}
\setlist[itemize]{leftmargin=1.35em,itemsep=0.2em,topsep=0.25em}
\setlist[enumerate]{leftmargin=1.55em,itemsep=0.35em,topsep=0.25em}
\newcolumntype{L}{>{\raggedright\arraybackslash}X}

% Keep section labels plain and compact so the guide remains easy to scan.
\makeatletter
\renewcommand\section{\@startsection{section}{1}{\z@}%
  {0.9em}% space before
  {0.25em}% space after
  {\normalfont\large\bfseries}}
\makeatother

\pagestyle{fancy}
\fancyhf{}
\renewcommand{\headrulewidth}{0pt}
\fancyfoot[L]{\scriptsize \CourseCodes\quad |\quad AI Fluency quiz guide}
\fancyfoot[R]{\scriptsize \thepage}

\newcommand{\GuideHeader}{%
  {\LARGE\bfseries AI Fluency Quizzes: Student Guide}\par
  \vspace{0.25em}
  \CourseCodes: \CourseName\hfill\CourseTerm\par
  Instructor: \InstructorName\hfill Virginia Tech
  \vspace{0.4em}\hrule\vspace{0.55em}
}

\begin{document}
\GuideHeader

\section*{What this quiz format is}
An AI Fluency quiz asks you to evaluate a short AI-generated answer about the technical material in the course. The answer will usually mix good observations with a claim that is incomplete, unsupported, or too confident. Your job is not to declare that ``AI is bad.'' Your job is to decide exactly what the evidence supports and what it does not.

The posted sample quiz is \textbf{not for credit}. It shows the structure we may use on future in-class quizzes.

\section*{How the two-pass quiz works}
\begin{enumerate}
  \item \textbf{Cold audit.} Read only the front. Mark one strong claim, one claim you would verify, and the most consequential leap. Record an initial disposition.
  \item \textbf{Course instruction.} During the lecture, learn the relevant technical ideas. Keep your original marks visible.
  \item \textbf{Evidence audit.} Turn the sheet over when instructed. Use the new evidence to audit one claim, complete a short technical check, and request a useful verification artifact.
  \item \textbf{Revision.} Give a final disposition. You may change your mind, or keep the same verdict for a better reason. Either can demonstrate learning.
\end{enumerate}

\section*{A useful answer structure}
For the main audit, move through four questions:

\begin{tabularx}{\linewidth}{@{}p{0.95in}L@{}}
  \toprule
  \textbf{Claim} & What, exactly, is the AI asserting? Restate it in your own words. \\
  \textbf{Evidence} & Which number, assumption, definition, or missing test bears on that claim? \\
  \textbf{Action} & Does the evidence justify only a comparison, a threshold choice, a pilot, or a real deployment? \\
  \textbf{Next step} & What concrete artifact would reduce uncertainty, and what result would make you stop? \\
  \bottomrule
\end{tabularx}

\section*{What the quiz assesses}
\begin{itemize}
  \item \textbf{Foundational machine learning:} definitions, assumptions, confusion-matrix rates, threshold tradeoffs, calibration, conditional probability, validation, and data scope.
  \item \textbf{Technical judgment:} choosing the right denominator, separating a model score from a real-world guarantee, and recognizing when evidence is too narrow for the proposed action.
  \item \textbf{AI fluency:} tracing claims to evidence, noticing confident but unsupported transitions, asking for falsifiable checks, and retaining responsibility for the final decision.
\end{itemize}

\section*{Practical expectations}
\begin{itemize}
  \item Bring something to write with; a second ink color is helpful but optional. Turn the page only when instructed.
  \item AI tools may not be used during an in-class quiz unless the instructor explicitly authorizes them or an approved accommodation applies.
  \item Treat the scenario as a technical decision exercise. The mushroom example is fictional and must never be used for real identification.
\end{itemize}

\newpage
\GuideHeader

\section*{What strong work looks like}
A strong response is specific and economical. It usually does four things:

\begin{itemize}
  \item names a particular claim rather than criticizing the entire answer;
  \item cites packet evidence or completes the relevant calculation correctly;
  \item explains why that evidence changes the proposed action; and
  \item asks for a test or artifact with a clear stop condition.
\end{itemize}

You do \textbf{not} need to diagnose every sentence. Prioritize the flaw that could cause the most consequential decision error.

\section*{Mini-example}
Suppose an AI writes: ``The classifier is 95\% accurate, so it is ready to deploy.'' A useful audit would not merely say ``unsupported.'' It would explain that overall accuracy hides which class was harmed, request confusion counts at the proposed operating threshold, and state what high-cost error would block deployment.

That response connects a foundational concept to an action. It also makes the next AI interaction more productive because the follow-up request is concrete and testable.

\section*{How we grade}
The main focus is your \textbf{thinking process}: whether you identify a consequential claim, connect it to evidence, explain what that evidence authorizes, and make a defensible decision. Open-ended judgment questions may have more than one strong answer. You do not need to match a preferred conclusion if your reasoning is technically sound and well supported.

When a question has a clear right or wrong answer---for example, a calculation, conditional probability, or definition---the answer and supporting work are graded for correctness.

The quiz follows the course's short in-class 2 / 1 / 0 policy:

\begin{itemize}
  \item \textbf{2 points:} clear, evidence-linked reasoning and correct work when a question has a single answer;
  \item \textbf{1 point:} relevant progress with a missing reasoning link or a technical error; and
  \item \textbf{0 points:} absent or materially blank.
\end{itemize}

Your initial cold judgment is not graded for being ``right.'' What matters is making your reasoning visible and improving it when the evidence warrants revision.

\section*{Why the quiz is designed this way}
AI answers can be fluent even when the evidence-to-conclusion step is weak. The two-pass design captures your first reaction, then lets you use the lecture's technical ideas to reassess it. This checks understanding more directly than recognizing an option in a multiple-choice list, and it treats revision as evidence of learning rather than as a mistake.

The paper format also matters: your original marks, later evidence, and final decision remain visible together. That creates a small record of how your reasoning changed.

\textbf{Best preparation:} understand the lecture's core concepts well enough to ask, ``What evidence would have to be true before I could rely on this conclusion?''

\end{document}
